\documentclass{ximera}

\input{../../preamble.tex}

\title{Function Analysis}

\begin{document}

\begin{abstract}
%
\end{abstract}
\maketitle






What do we want when analyzing a function? 




We would like to analyze all of their characteristics and features. 



\begin{itemize}
     \item \textbf{\textcolor{red!70!black}{Domain}} 
     \item \textbf{\textcolor{red!70!black}{Zeros}} 
     \item \textbf{\textcolor{red!70!black}{Continuity}} \\
     \text{    }$\circ$ \textbf{\textcolor{purple!85!blue}{discontinuities}}\\
     \text{    }$\circ$ \textbf{\textcolor{purple!85!blue}{singularities}}\\
     \item \textbf{\textcolor{red!70!black}{End-Behavior}} 
     \item \textbf{\textcolor{red!70!black}{Behavior}} \\
     \text{    }$\circ$ \textbf{\textcolor{purple!85!blue}{intervals where increasing}}\\
     \text{    }$\circ$ \textbf{\textcolor{purple!85!blue}{intervals where decreasing}}\\
     \item \textbf{\textcolor{red!70!black}{Global Maximum and Minimum}} 
     \item \textbf{\textcolor{red!70!black}{Local Maximums and Minimums}} 
     \item \textbf{\textcolor{red!70!black}{Range}} 
     \item \textbf{\textcolor{blue!55!black}{...and we would like a nice graph}} 
\end{itemize}






In addition, we want \textbf{\textcolor{red!70!black}{algebraic}} and functional reasoning on how we arrive at our conclusions.


To help us with this goal of algebraic communication, we can use a nice graph.  Everything we eventually want to say algebraically corresponds to graphical characteristics.

We don't want our reasoning to be ``look at the graph'', however, we can look at the graph and have it help us translate our reasoning into algebra and function language.





Along with dots representing function pairs, a nice graph, including auxillary graphing items
\begin{itemize}
	\item intercepts
	\item endpoints
	\item vertical asymptotes
	\item horizontal asymptotes
\end{itemize}





A graph is a communication tool.




\subsection*{Learning Outcomes}

\begin{sectionOutcomes}
In this section, students will 

\begin{itemize}
\item analyze functions from their formula.
\end{itemize}
\end{sectionOutcomes}










\begin{center}
\textbf{\textcolor{green!50!black}{ooooo-=-=-=-ooOoo-=-=-=-ooooo}} \\

more examples can be found by following this link\\ \link[More Examples of Function Analysis]{https://ximera.osu.edu/csccmathematics/precalculus2/precalculus2/functionAnalysis/examples/exampleList}

\end{center}









\end{document}
